%========================================================
% APPENDIX C.3 — Statistical Composite Fiber-Channel Analytics
%========================================================

\section{Statistical Composite Fiber-Channel Analytics}
\label{app:statistical_composite_channel_analytics}

This appendix derives the analytical expressions associated with the statistical composite fiber model introduced in Section~\ref{sec:statistical_composite_fiber_channel}.

\medskip

The model combines:

\begin{itemize}

\item statistical phase averaging,

\item local isotropic depolarization acting independently on the two propagation arms.

\end{itemize}

\medskip

The reference Bell state is

\[
|\Psi^+\rangle
=
\frac{1}{\sqrt{2}}
\left(
|01\rangle
+
|10\rangle
\right),
\]

with density operator

\[
\rho_0
=
|\Psi^+\rangle
\langle
\Psi^+|.
\]

The phase-evolved pure state is

\begin{equation}
|\psi(\theta)\rangle
=
\frac{1}{\sqrt{2}}
\left(
|01\rangle
+
e^{i\theta}|10\rangle
\right).
\label{eq:appendix_statistical_phase_state}
\end{equation}

\medskip

Define the phase-coherence parameter

\begin{equation}
\mu
=
\left\langle
e^{i\theta}
\right\rangle
=
\sum_k
p_k
e^{i\theta_k},
\label{eq:appendix_mu_definition}
\end{equation}

and the depolarization contraction factor

\begin{equation}
\eta
=
(1-p_A)(1-p_B).
\label{eq:appendix_eta_definition}
\end{equation}

%--------------------------------------------------------
% Ensemble density operator
%--------------------------------------------------------

\subsection*{Ensemble density operator}

The ensemble-averaged density operator is

\begin{equation}
\rho^{(\mathrm{ens})}
=
\sum_k
p_k
|\psi(\theta_k)\rangle
\langle
\psi(\theta_k)|.
\label{eq:appendix_ensemble_density_operator}
\end{equation}

Using Eq.~\eqref{eq:appendix_statistical_phase_state}, one obtains

\begin{equation}
\rho^{(\mathrm{ens})}
=
\frac12
\begin{pmatrix}
0 & 0 & 0 & 0
\\
0 & 1 & \mu^* & 0
\\
0 & \mu & 1 & 0
\\
0 & 0 & 0 & 0
\end{pmatrix}.
\label{eq:appendix_ensemble_density_matrix}
\end{equation}

\medskip

Including local isotropic depolarization contracts the bipartite correlations by the factor \( \eta \). The resulting effective state is

\begin{equation}
\rho_{\eta}^{(\mathrm{ens})}
=
\frac14
\left[
I\otimes I
+
\sum_{i,j}
T_{ij}^{(\mathrm{ens})}
\sigma_i\otimes\sigma_j
\right],
\label{eq:appendix_composite_ensemble_state}
\end{equation}

where the tensor components are derived below.

%--------------------------------------------------------
% Correlation tensor
%--------------------------------------------------------

\subsection*{Correlation tensor}

For the ensemble state without local depolarization, the correlation tensor is

\begin{equation}
T^{(\mathrm{ens})}
=
\begin{pmatrix}
\mathrm{Re}(\mu)
&
-\mathrm{Im}(\mu)
&
0
\\
\mathrm{Im}(\mu)
&
\mathrm{Re}(\mu)
&
0
\\
0
&
0
&
-1
\end{pmatrix}.
\label{eq:appendix_ensemble_tensor}
\end{equation}

\medskip

Under local isotropic depolarization,

\begin{equation}
T_{\eta}^{(\mathrm{ens})}
=
\eta
T^{(\mathrm{ens})},
\label{eq:appendix_ensemble_tensor_contraction}
\end{equation}

thus

\begin{equation}
T_{\eta}^{(\mathrm{ens})}
=
\eta
\begin{pmatrix}
\mathrm{Re}(\mu)
&
-\mathrm{Im}(\mu)
&
0
\\
\mathrm{Im}(\mu)
&
\mathrm{Re}(\mu)
&
0
\\
0
&
0
&
-1
\end{pmatrix}.
\label{eq:appendix_final_ensemble_tensor}
\end{equation}

The tensor components are therefore

\begin{align}
T_{xx}
&=
\eta\,\mathrm{Re}(\mu),
\\
T_{xy}
&=
-\eta\,\mathrm{Im}(\mu),
\\
T_{yx}
&=
\eta\,\mathrm{Im}(\mu),
\\
T_{yy}
&=
\eta\,\mathrm{Re}(\mu),
\\
T_{zz}
&=
-\eta.
\end{align}

All remaining tensor components vanish.

%--------------------------------------------------------
% Global purity
%--------------------------------------------------------

\subsection*{Global purity}

The global purity is

\begin{equation}
\gamma
\left(
\rho_{AB}^{(\mathrm{ens})}
\right)
=
\operatorname{Tr}
\left[
\left(
\rho_{AB}^{(\mathrm{ens})}
\right)^2
\right].
\end{equation}

Using the Pauli-basis representation,

\begin{equation}
\rho
=
\frac14
\left[
I\otimes I
+
\sum_{i,j}
T_{ij}
\sigma_i\otimes\sigma_j
\right],
\end{equation}

together with Pauli orthogonality,

\begin{equation}
\operatorname{Tr}
\left[
(\sigma_i\otimes\sigma_j)
(\sigma_k\otimes\sigma_l)
\right]
=
4
\delta_{ik}
\delta_{jl},
\end{equation}

one obtains

\begin{equation}
\gamma
=
\frac14
\left[
1
+
\sum_{i,j}
T_{ij}^2
\right].
\label{eq:appendix_purity_tensor_formula}
\end{equation}

Substituting Eq.~\eqref{eq:appendix_final_ensemble_tensor},

\begin{align}
\sum_{i,j}
T_{ij}^2
&=
\eta^2
\left[
2|\mu|^2
+
1
\right].
\end{align}

Therefore,

\begin{equation}
\boxed{
\gamma
\left(
\rho_{AB}^{(\mathrm{ens})}
\right)
=
\frac{
1+\eta^2(1+2|\mu|^2)
}{4}
}
\label{eq:appendix_ensemble_purity}
\end{equation}

\medskip

For \( \eta=1 \), one recovers

\begin{equation}
\gamma
\left(
\rho_{AB}^{(\mathrm{ens})}
\right)
=
\frac{1+|\mu|^2}{2}.
\label{eq:appendix_phase_ensemble_purity}
\end{equation}

%--------------------------------------------------------
% Reduced density operators
%--------------------------------------------------------

\subsection*{Reduced density operators}

The reduced density operators remain maximally mixed:

\begin{equation}
\rho_A
=
\rho_B
=
\frac{I_2}{2}.
\end{equation}

Therefore,

\begin{equation}
\boxed{
\gamma_A
=
\gamma_B
=
\frac12
}.
\label{eq:appendix_ensemble_reduced_purity}
\end{equation}

%--------------------------------------------------------
% Fidelity
%--------------------------------------------------------

\subsection*{Fidelity with respect to \texorpdfstring{\( |\Psi^+\rangle \)}{|Psi+>}}

The Bell-state fidelity is

\begin{equation}
F_{\Psi^+}
=
\langle
\Psi^+
|
\rho_{\eta}^{(\mathrm{ens})}
|
\Psi^+
\rangle.
\end{equation}

Using the ensemble density operator,

\begin{equation}
F_{\Psi^+}
=
\frac12
\left[
1
+
\eta\,\mathrm{Re}(\mu)
\right].
\end{equation}

Therefore,

\begin{equation}
\boxed{
F_{\Psi^+}
=
\frac{
1+\eta\,\mathrm{Re}(\mu)
}{2}
}
\label{eq:appendix_ensemble_fidelity}
\end{equation}

\medskip

For \( \eta=1 \), this reduces to

\begin{equation}
F_{\Psi^+}
=
\frac{
1+\mathrm{Re}(\mu)
}{2}.
\end{equation}

%--------------------------------------------------------
% Concurrence
%--------------------------------------------------------

\subsection*{Concurrence}

The ensemble density operator belongs to the class of \(X\)-states. For a two-qubit \(X\)-state,

\begin{equation}
C
=
2
\max
\left(
0,
|\rho_{23}|
-
\sqrt{\rho_{11}\rho_{44}}
\right).
\label{eq:appendix_xstate_concurrence}
\end{equation}

For the present model,

\begin{equation}
\rho_{23}
=
\frac{\eta\mu}{2},
\end{equation}

while the depolarizing contribution generates

\begin{equation}
\rho_{11}
=
\rho_{44}
=
\frac{1-\eta}{4}.
\end{equation}

Therefore,

\begin{align}
C
&=
2
\max
\left[
0,
\frac{\eta|\mu|}{2}
-
\frac{1-\eta}{4}
\right]
\nonumber
\\
&=
\max
\left[
0,
\eta|\mu|
-
\frac{1-\eta}{2}
\right].
\end{align}

Hence,

\begin{equation}
\boxed{
C
=
\max
\left[
0,
\eta|\mu|
-
\frac{1-\eta}{2}
\right]
}
\label{eq:appendix_ensemble_concurrence}
\end{equation}

%--------------------------------------------------------
% CHSH nonlocality
%--------------------------------------------------------

\subsection*{CHSH nonlocality}

The maximal CHSH parameter is determined by the Horodecki criterion:

\begin{equation}
S_{\max}
=
2
\sqrt{
\lambda_1+\lambda_2
},
\end{equation}

where \( \lambda_1,\lambda_2 \) are the two largest eigenvalues of

\[
T^TT.
\]

\medskip

Using Eq.~\eqref{eq:appendix_final_ensemble_tensor},

\begin{equation}
T^TT
=
\eta^2
\begin{pmatrix}
|\mu|^2 & 0 & 0
\\
0 & |\mu|^2 & 0
\\
0 & 0 & 1
\end{pmatrix}.
\end{equation}

The two largest eigenvalues are therefore

\[
\lambda_1
=
\eta^2,
\qquad
\lambda_2
=
\eta^2|\mu|^2.
\]

Thus,

\begin{equation}
\boxed{
S_{\max}
=
2\eta
\sqrt{
1+|\mu|^2
}
}
\label{eq:appendix_ensemble_chsh}
\end{equation}

Bell-inequality violation requires

\begin{equation}
S_{\max}>2,
\end{equation}

which gives

\begin{equation}
\boxed{
\eta
\sqrt{
1+|\mu|^2
}
>
1
}
\label{eq:appendix_ensemble_chsh_threshold}
\end{equation}

%--------------------------------------------------------
% Summary
%--------------------------------------------------------

\subsection*{Summary}

Collecting the analytical expressions:

\[
\boxed{
\begin{aligned}
T_{xx}
&=
\eta\,\mathrm{Re}(\mu),
\\
T_{xy}
&=
-\eta\,\mathrm{Im}(\mu),
\\
T_{yx}
&=
\eta\,\mathrm{Im}(\mu),
\\
T_{yy}
&=
\eta\,\mathrm{Re}(\mu),
\\
T_{zz}
&=
-\eta,
\\[2mm]
\gamma
\left(
\rho_{AB}^{(\mathrm{ens})}
\right)
&=
\frac{
1+\eta^2(1+2|\mu|^2)
}{4},
\\
\gamma_A
&=
\gamma_B
=
\frac12,
\\
F_{\Psi^+}
&=
\frac{
1+\eta\,\mathrm{Re}(\mu)
}{2},
\\
C
&=
\max
\left[
0,
\eta|\mu|
-
\frac{1-\eta}{2}
\right],
\\
S_{\max}
&=
2\eta
\sqrt{
1+|\mu|^2
}.
\end{aligned}
}
\]

These expressions provide the analytical reference for the statistical composite fiber-channel model introduced in Section~\ref{sec:statistical_composite_fiber_channel}.